% PIBA Template -- Software Architecture Lab, WS 2026/27, University of Stuttgart
%
% Self-contained: compile with pdflatex (no additional files needed).
% Rules: the whole PIBA must fit on ONE page. Do not change font size or margins.
% Replace all text in \todo{...} and remove the hints before submitting.

\documentclass[11pt,a4paper]{article}

\usepackage[utf8]{inputenc}
\usepackage[T1]{fontenc}
\usepackage{lmodern}
\usepackage[a4paper,margin=2cm]{geometry}
\usepackage{enumitem}
\usepackage{graphicx}
\usepackage{xcolor}
\usepackage[hidelinks]{hyperref}

\pagestyle{empty}
\setlength{\parindent}{0pt}
\setlength{\parskip}{0.5em}
\setlist{itemsep=0.1em, topsep=0.2em}

% Placeholder marker; leave the definition, but remove all uses before submitting.
\newcommand{\todo}[1]{\textcolor{gray}{[#1]}}

% Part heading
\newcommand{\pibapart}[1]{\vspace{0.4em}{\large\textbf{#1}}\par}

% ---------------------------------------------------------------------------
% Project information
% ---------------------------------------------------------------------------
\newcommand{\projecttitle}{Dance Teacher}
\newcommand{\teamname}{Vienna}
\newcommand{\teammembers}{Aylina Ashkenov, Leon Layer, Stevan Paunovic, Nils Kwiatkowski, Sohan Amarapilli, Tobias Kostorz}

\begin{document}

\begin{center}
  {\LARGE\textbf{\projecttitle}}\\[0.4em]
  {\large PIBA -- Software Architecture Lab, WS 2026/27}\\[0.3em]
  \teamname\ \textbullet\ \teammembers
\end{center}

% ---------------------------------------------------------------------------
\pibapart{Problem}
% What is wrong, missing, or hard today, and for whom?
% Describe the situation today, not your solution. Name the system, the users,
% and the concrete pain. Quantify where possible.
Learning to dance without professional guidance can be challenging. Learners
practising dance routines with only reference videos must identify movement and
timing errors themselves. Video playback alone does not provide individual
corrections. A camera-based assistant can distinguish actual deviations from
unreliable tracking and deliver feedback before it becomes outdated. Our target
setting is individual practice of short routines with a fixed camera and the
learner's full body visible.

% ---------------------------------------------------------------------------
\pibapart{Idea}
% What is your proposed solution, and what does its architecture look like?
% The architecture is the heart of your idea.
Dance Teacher overlays a skeleton on the learner's camera image. It compares
the learner's poses with a reference video that the user can upload and gives
feedback. Movement similarities are displayed. Everything runs locally on one
laptop.

The system has two pipelines that use the same pose estimation and only share a
file with the reference poses, so each can be changed on its own. The reference
pipeline runs once per video and marks key poses. The practice loop runs for
every camera frame and compares the learner's normalized pose with the key pose
the learner is currently practising. Teaching methods are interchangeable
learning strategies (Strategy pattern). The first one is Pose by Pose; more can
be added later without changing other components.

The practice loop is a chain of steps, each passing its result to the next
(Pipes and Filters). The pose library and the camera sit behind our own
interfaces, so we can replace them (Ports and Adapters). Key quality
attributes: fast feedback, robustness against tracking errors and easy
extension.

Out of scope are model training, partner dancing, multiple cameras and detailed
evaluation of dance technique.

% ---------------------------------------------------------------------------
\pibapart{Benefits}
% What improves if the idea works, and how would you show it?
% Link each benefit back to the problem.
\begin{itemize}
  \item \textbf{Improvement:} Better dance moves and the possibility to teach
  yourself with actionable feedback -- supports self-correction.
  \textbf{Shown by:} a field study comparing learning a dance from a video with
  vs.\ without the application.
  \item \textbf{Improvement:} Feedback the learner can trust -- tracking noise
  should not lead to false corrections, and feedback should arrive in time.
  \textbf{Shown by:} tests with recorded videos.
\end{itemize}

% ---------------------------------------------------------------------------
\pibapart{Action}
% What will you build and evaluate, in which steps, within 3 months?
% Mark the minimum deliverable and keep some buffer.
\begin{enumerate}
  \item Nov 2026 -- Requirements specification + domain model
  \item Dec 2026 -- Architecture; first tests with the pose-estimation library
  \item Jan 2027 -- \textbf{Minimum deliverable:} a first end-to-end prototype
  with one learning strategy, Pose by Pose: the learner takes the key poses of
  the reference video one after another, and the next pose only appears once
  the current one is matched.
  \item Feb 2027 -- Evaluation and presentation; if time allows, more learning strategies.
\end{enumerate}
\textbf{Required technologies, data, and infrastructure:} a laptop with a
webcam, a pretrained pose-estimation library, and reference videos with a
single dancer.\\
\textbf{Team split:} dynamic split between the members Tobias, Leon, Aylina,
Nils, Sohan and Stevan.

% ---------------------------------------------------------------------------
\vfill
{\small\textbf{AI usage statement.} ChatGPT and Claude were used for
brainstorming, rephrasing and improving the text of the project description.}

\end{document}
